\section{Introduction}
Actor insertion pipelines can produce plausible point sets while violating the sensing
process that generated them. Evaluation becomes especially misleading when held-out
returns influence the forward pass, output density differs across methods, or a learned
readout hides weak geometry. A useful simulator must distinguish invalid and censored
queries, valid no-returns, background occlusion, and Actor returns while composing all
surfaces in world-depth order.

We organize the study around that task rather than around an inherited model. Our first
route asks whether mature point completion gains physical consistency from measured
hit and free-space constraints. Our second asks whether the broader problem instead
requires a complete neural LiDAR model. Both face deterministic baselines and external
methods before any final test is opened.

This audit contributes (i) a target-free, ray-preserving Actor data contract with
log-level role controls; (ii) a matched geometry and conditional-return evaluation that
exposes density and readout effects; (iii) executed AdaPoinTr and LiDAR4D baselines; and
(iv) a pre-registered negative route decision on 501 clean-development Actors. The
negative result is central: the proposed observation-constrained completion does not
beat Actor-local TSDF, so we stop before route selection, source testing, and application
claims.
